% Method. Contract: ../references/section-playbook.md and equation-grammar.md
% 1300-2100 words, 3-4 subsections, at least 4 numbered equations.
% The overview figure is the first thing in the section, before any equation.
\section{Method}
\label{sec:method}
% Tier C structural overview. First float of the Method section, placed before
% the section's first equation. Carries no magnitudes: it states what is fixed,
% what varies, and what the paper reports.
\begin{figure*}[t]
  \centering
  \resizebox{0.99\textwidth}{!}{%
  \begin{tikzpicture}[x=1cm,y=1cm]
    \node[acbox,minimum width=3.5cm,minimum height=1.15cm,align=center] (a) at (1.9,1.1)
      {\textbf{Frozen \evalsem}\\Same bound code\\Same evaluation precision};
    \node[acbox,minimum width=3.5cm,minimum height=1.15cm,align=center] (b) at (6.7,1.1)
      {\textbf{Independent factors}\\\splitseed\\\trajseed};
    \node[acbox,minimum width=3.7cm,minimum height=1.15cm,align=center] (c) at (11.7,1.1)
      {\textbf{Comparable outputs}\\Mean, standard deviation,\\range, paired contrasts};
    \draw[acarrow] (a) -- (b);
    \draw[acarrow] (b) -- (c);
    % text width is a cap, not a floor: at this 4.1cm center spacing the two
    % callouts can grow only downward (more lines), never sideways into each
    % other, regardless of how long the caption text is.
    \node[acsoft,text width=3.6cm,align=center] at (5.0,2.55)
      {The selection method is the treatment, not a source of hidden protocol
       drift};
    \node[acsoft,text width=3.6cm,align=center] at (9.1,2.55)
      {Every reported statistic traces to preserved primary runs};
  \end{tikzpicture}}
  \caption{\textbf{Study design.} The evaluation is fixed before any comparison
    begins, the two sources of randomness are varied independently, and the
    reported statistics are defined in advance. Fixing the left panel is what
    makes the right panel a measurement of the method rather than of the
    protocol.}
  \label{fig:overview}
\end{figure*}


\subsection{Setting and notation}
\label{sec:setting}

A disagreement-based certificate bounds the risk of a target model through a
simpler \surrogate\ that admits a computable guarantee.
\Cref{fig:overview} summarizes the design: the evaluation is fixed first, the
two sources of randomness are varied independently, and the reported statistics
are defined before any comparison is run. We write one evaluation instance as a
tuple $(\mathcal{S}, m, \tau) \rightarrow \bnd$, where $\mathcal{S}$ is the data
split, $m$ the selection method that determines which examples the \surrogate\
is built from, $\tau$ the \trajseed\ that fixes the optimization path, and
$\bnd$ the resulting \certbound. Lower $\bnd$ is better.

The certificate has two parts that move in opposite directions when the
selection method changes, which is the reason a selection effect is not
self-evident:
\begin{equation}\label{eq:certificate}
  \bnd \;=\; \underbrace{\mathcal{C}(\hypclass)}_{\text{\surrogate\ guarantee}}
  \;+\;
  \underbrace{\widehat{d}(\dataset_u)}_{\text{\disagreement}},
\end{equation}
where $\mathcal{C}(\hypclass)$ is the \surrogate's own guarantee, which depends
on its complexity relative to its prior, and $\widehat{d}(\dataset_u)$ is the
estimated disagreement between target and \surrogate\ measured on an unlabeled
set $\dataset_u$. \Cref{eq:certificate} is what makes the comparison in this
paper non-trivial: a selection method that reduces the second term by fitting
$\dataset_u$ more aggressively can enlarge the first by more than it gains, and
the observed $\bnd$ reports only the sum. Neither term is separately reported by
the pipeline, so a method comparison at the level of $\bnd$ cannot attribute a
difference to one term or the other without additional instrumentation.

Three things are fixed for every run reported in this paper: the bound
computation itself, the evaluation precision, and the split boundaries between
training, validation, disagreement, and test data. Changing any of them makes
numbers on either side of the change incomparable, so they are frozen before the
matrix is registered and are never treated as tunable. Concretely, the
confidence parameter is $\delta = 0.05$, the bound-inversion grid holds 1000
points, the Monte Carlo estimate of the \disagreement uses 10{,}000 samples,
and the four-way split is drawn once per \splitseed\ and reused by every method
at that split. The training and disagreement partitions never overlap.

\subsection{The compared selection methods}
\label{sec:methods}

The two selection methods differ in exactly one respect: the criterion used to
choose the examples the \surrogate\ is fitted on. \methodrandom\ draws them
uniformly without replacement from the training partition. \methodmax\ ranks
candidates by the target model's loss and takes the highest-loss examples, which
is the strategy with the most direct control over the region where target and
\surrogate\ are most likely to disagree. Both draw the same number of examples,
run for the same number of epochs, use the same optimizer and schedule, and are
evaluated by the same frozen code path. The \methodcontrol\ is the unmodified
upstream pipeline, re-run and recorded in this environment rather than quoted
from a prior publication, because published numbers came from different
hardware, a different epoch budget, and full-precision evaluation.

Holding the budget equal needs saying precisely, because the two criteria
consume it differently. Both arms select the same number of examples and train
for the same number of epochs on them, so the \surrogate\ sees the same number
of gradient steps in both. \methodmax\ additionally requires a forward pass over
the training partition to rank candidates by target loss, which is a one-time
cost outside the training loop and is not charged against the epoch budget. That
asymmetry favours \methodmax\ in wall-clock terms and is irrelevant to the
certificate, which depends on the fitted \surrogate\ and not on how its training
set was chosen. We report it because a budget claim that is only true of one
resource is the kind of detail that makes two studies incomparable later.

This is a deliberately narrow contrast. Richer strategies exist, including
disagreement-aware and hybrid criteria, and they are not evaluated here. The
narrow contrast is what makes the identification argument in the next subsection
possible: two arms that differ in one criterion, run over the same splits with
the same budget, admit an exact pairing that a broader sweep would not.

\subsection{Separating the factors}
\label{sec:factors}

A comparison between two selection methods is identifiable only if the remaining
sources of variation are either held fixed or averaged over. We therefore treat
each observed bound as the sum of a split effect and a trajectory residual under
that split:
\begin{equation}\label{eq:decomposition}
  \bnd_{\splitidx\trajidx}^{(m)}
  \;=\;
  \mu^{(m)} + \alpha_{\splitidx} + \varepsilon_{\splitidx\trajidx}^{(m)},
  \qquad
  \textstyle\sum_{\splitidx} \alpha_{\splitidx} = 0,
\end{equation}
where $\bnd_{\splitidx\trajidx}^{(m)}$ is the bound obtained by method $m$ at
\splitseed\ $\splitidx$ and \trajseed\ $\trajidx$, $\mu^{(m)}$ is the method's
grand mean, $\alpha_{\splitidx}$ is the effect of the split itself, shared by all
methods, and $\varepsilon_{\splitidx\trajidx}^{(m)}$ is the trajectory residual.
\Cref{eq:decomposition} states the identification problem directly: a difference
between methods measured at different splits confounds $\mu^{(m)}$ with
$\alpha_{\splitidx}$, so it is not evidence about the method.

\textbf{Design rule}: \textit{a method contrast is reported only when it is
measured within a split, and a method ranking is reported only when it is
consistent across splits.}

Both halves of that rule matter. The first removes $\alpha_{\splitidx}$ from the
contrast by construction rather than by assumption. The second prevents a
ranking that holds at one split and reverses at another from being reported as a
property of the method, which is the failure mode that a single-split comparison
cannot detect. The rule is restrictive on purpose: it can leave a study with a
measured contrast and no ranking to report, and that outcome is a result rather
than a shortfall.

\subsection{Estimators and uncertainty}
\label{sec:estimators}

For each method $m$ and \splitseed\ $\splitidx$ we run $T$ trajectory seeds and
summarize them by their mean and their spread:
\begin{equation}\label{eq:persplit}
  \bar{\bnd}^{(m)}_{\splitidx} = \frac{1}{T}\sum_{\trajidx=1}^{T}
    \bnd^{(m)}_{\splitidx\trajidx},
  \qquad
  s^{(m)}_{\splitidx} = \sqrt{\frac{1}{T-1}\sum_{\trajidx=1}^{T}
    \bigl(\bnd^{(m)}_{\splitidx\trajidx} - \bar{\bnd}^{(m)}_{\splitidx}\bigr)^2},
\end{equation}
where $T$ is the number of trajectory seeds per split,
$\bar{\bnd}^{(m)}_{\splitidx}$ the per-split mean, and $s^{(m)}_{\splitidx}$ the
trajectory standard deviation under a fixed split. \Cref{eq:persplit} is the
unit of evidence in this paper: a single trajectory is an observation, and only
$\bar{\bnd}^{(m)}_{\splitidx}$ with its spread is quoted as a result. The spread
is reported alongside the mean everywhere the mean appears, including in the
abstract, because the two numbers answer different questions and the first is
uninterpretable without the second.

The headline comparison is the paired contrast between two methods, taken within
split and trajectory index so that both $\alpha_{\splitidx}$ and any shared
trajectory effect cancel:
\begin{equation}\label{eq:contrast}
  \Delta_{\splitidx} = \bar{\bnd}^{(m_1)}_{\splitidx} - \bar{\bnd}^{(m_2)}_{\splitidx},
  \qquad
  \bar{\Delta} = \frac{1}{S}\sum_{\splitidx=1}^{S} \Delta_{\splitidx},
\end{equation}
where $S$ is the number of \splitseed{}s, $\Delta_{\splitidx}$ the within-split
contrast, and $\bar{\Delta}$ its mean across splits. \Cref{eq:contrast} is the
quantity reported in the main results table; the per-split $\Delta_{\splitidx}$
are reported alongside it so that a mean contrast cannot hide a reversal.

Finally, a contrast is only reported as a finding when it exceeds the noise it is
measured against:
\begin{equation}\label{eq:decision}
  \bigl|\bar{\Delta}\bigr| \;>\; \kappa \cdot
  \max_{m,\splitidx}\, s^{(m)}_{\splitidx},
  \qquad \kappa = 1.0,
\end{equation}
where $\kappa$ is a threshold fixed before the matrix was run and
$\max_{m,\splitidx} s^{(m)}_{\splitidx}$ is the largest within-split trajectory
spread observed anywhere in the matrix. The threshold is expressed against an observed spread rather than against a
distributional assumption because the quantity in question, the movement of a
certificate under a change of \trajseed, has no established parametric form here
and only three observations per cell from which to estimate one.
\Cref{eq:decision} is a stopping rule rather than a significance test: when it fails, the finding this paper reports is
that the methods are indistinguishable at this scale, and that statement is
recorded rather than replaced by a smaller comparison that happens to pass. The
comparison against the largest observed spread, rather than against the spread of
the winning arm, is the conservative choice; \cref{sec:appendix-statistics}
reports what the verdict would have been under the alternatives.

\subsection{Provenance}
\label{sec:provenance}

Every run writes a raw result together with a manifest recording its method,
source script and revision, exact command, dataset, \splitseed, \trajseed, full
configuration, timestamps, and environment. Aggregates enumerate the primary run
paths they were computed from. A quantity whose primary record is missing,
overwritten, or method-ambiguous is excluded from the main paper rather than
reported with a caveat. This matters in practice because the upstream pipeline
names its output files from hyperparameters alone, so two different methods run
at the same hyperparameters write to the same filename and the second silently
overwrites the first. The aggregation step reads only manifests and raw results. It refuses to write
an aggregate when a listed primary input is absent from disk, so a broken
evidence chain surfaces as a failure at aggregation time rather than as a
plausible number in a table. Each aggregate stores the estimator it applied, the
per-seed values it combined, and the paths it combined them from, which is what
makes a reported statistic checkable rather than merely reproducible in
principle. \Cref{sec:appendix-provenance} gives the manifest schema and the
aggregation procedure in full.
